\subsection{False-positive audit}
\label{sec:human_null_audit}
The 128-source confirmation contains only 128 human-null prefixes, too few to assess a nominal $0.1\%$ error rate precisely. We therefore freeze a separate 4,096-prefix diagnostic before computing its owner scores. It uses HC3 human-labelled answers from question-hash buckets 65--89, outside the probability-interface evaluation and public-calibration buckets. All retained, development, confirmation and calibration question identities are excluded. Questions and first-100-word human prefixes are deduplicated, and known human-prefix hashes are excluded. No threshold or key is fitted on this set.

The native TextSeal2 v2 decision remains $p\leq0.001$, with null-tail $p$-value $p\in[0,1]$. It flags 5 of 4,096 prefixes: an empirical false-positive rate of $0.122\%$, with an exact binomial 95\% interval $[0.040\%,0.285\%]$. Every prefix has at least one scored context/token pair. This cohort is compatible with the nominal rate but does not precisely establish that rate for every key, dataset or adaptively selected input. The interval is descriptive under the usual independent-sample model, conditional on the existing key and this human-labelled population.

The audit runs entirely on CPU. It calls the released native copula-score primitive in batches, preserves the native first-occurrence context/token-pair rule, and uses the unchanged Gamma-tail function. Before computing the audit, it checks agreement with the scalar detector on 12 probes: empty, short and repeated-token cases plus eight fixed-position real prefixes. All scored-pair counts and per-token values agree to the declared $10^{-12}$ tolerance. The frozen manifest, individual $p$-values, source hash and result hashes remain in the audit receipt. This is an implementation and null-behavior check, not evidence of attack resistance or a new calibration procedure.
